\begin{finding}[\fid{F24}: signed Euler spectra at even path lengths]\label{f:signed}
Keep the pointwise setup and path identity of \fid{F03}, and put $w=a+d$.
Replace its input sign hypothesis by
\[
 M_n=\sum_{\alpha\in A}c_\alpha\alpha^n\quad(n\ge1),
 \qquad c_\alpha\in\mathbb R\setminus\{0\},
\]
where $A$ is finite, consists of distinct nonzero roots, and is nonempty.
Thus $A$ is the \emph{visible spectrum}, after grouping equal roots and deleting
zero coefficients. Suppose that for an unbounded set of \emph{even} $k\ge2$,
\[
 q^nN_{k,n}=-\sum_{\beta\in B_k}b_{k,\beta}\beta^n+E_{k,n},
 \qquad b_{k,\beta}>0,
\]
for every $n\ge1$, with the error bound of \fid{F04}. Then every
$\alpha\in A$ satisfies $|\alpha|\le q^{w/2}$. The positive coefficients
$b_{k,\beta}$ may be real; positive integer multiplicities suffice.

The all-embedding conclusion holds when each embedding has the corresponding
trace system, real transformed input coefficients, and error bounds. Integer
Euler multiplicities meet the coefficient condition. A reciprocal relation
$A=q^w/A$ among the visible supports, or a separately bounded visible dual
system as in \fid{F17}, supplies purity of the visible spectrum.

If the identities are available only at $n=e_kr$, the same conclusion holds
provided that the largest original radius remains visible after powering
by $e_k$, for an unbounded set of the allowed even $k$. This extra condition
cannot be omitted for signed spectra.
\end{finding}
\begin{proof}
Let $R=\max_{\alpha\in A}|\alpha|$ and
$A_R=\{\alpha:|\alpha|=R\}$. In the group algebra of the multiplicative group
generated by $A_R$, form
\[
 F=\sum_{\alpha\in A_R}c_\alpha[\alpha],\qquad
 F^k=\sum_\rho d_{k,\rho}[\rho],
\]
grouping equal products. The coefficients $d_{k,\rho}$ are real and
\[
 \sum_\rho d_{k,\rho}
   =\left(\sum_{\alpha\in A_R}c_\alpha\right)^k\ge0
\]
because $k$ is even. Moreover $F^k\ne0$. Otherwise evaluating
$[\alpha]\mapsto\alpha^n$ would give
$(\sum_{A_R}c_\alpha\alpha^n)^k=0$ for every $n\ge1$.
The Vandermonde matrix for $n=1,\ldots,|A_R|$ is invertible, contradicting
the nonzero coefficients on distinct nonzero roots. This argument uses
nonnilpotence of $F$, not a false assertion that a group algebra with
torsion has no zero divisors. Since a nonzero finite collection of
nonpositive real numbers has negative sum, at least one $d_{k,\rho}$ is
strictly positive.

Every root in $F^k$ has modulus $R^k$. A product containing a smaller-radius
root has smaller modulus, so it cannot cancel the chosen positive
coefficient. The path identity gives
\[
 M_n^k+\sum_{\beta\in B_k}b_{k,\beta}\beta^n=E_{k,n}-R_{k,n}.
\]
Adding the positive $B_k$ spectrum cannot cancel the chosen top-radius
coefficient either. The left side therefore has a surviving root of modulus
$R^k$. Its other coefficients need not be positive. Applying \fid{F02}
and \eqref{eq:overhead} for fixed $k$ gives
\[
 R^k\le q^{kw/2+\max(d/2+1,\eta_k)}.
\]
Take the $k$th root and the unbounded even-$k$ limit. The all-embedding and
dual arguments are the same as in \fid{F03} and \fid{F17}, but only visible
roots can be used in a reciprocal pairing.

For the finite-extension assertion, group first:
\[
 M_{er}=\sum_\gamma d_{e,\gamma}\gamma^r,
 \qquad d_{e,\gamma}=\sum_{\alpha^e=\gamma}c_\alpha.
\]
If the nonzero grouped spectrum is nonempty, put
$R_e=\max_{d_{e,\gamma}\ne0}|\gamma|^{1/e}$. Applying the preceding
argument in $r$ with $e=e_k$ gives
\[
 R_{e_k}\le q^{w/2+\max(d/2+1,\eta_k)/k}.
\]
Thus $R_{e_k}=R$ along an unbounded allowed set suffices. This equality is
automatic if the top-radius coefficients have one common sign, or if no
ratio of two distinct top-radius roots is a root of unity. It is not
automatic for general signed coefficients. Each embedding needs its own
top-radius visibility condition.
\end{proof}

\paragraph{Two exact limits of the refinement.}
First, take $a=0$, $d=1$, and let $t_n$ be the constant function $0$ for even
$n$ and $2$ for odd $n$. Then $|t_n|\le2$ and
\[
 M_n=q^n-(-q)^n.
\]
Both visible roots have modulus $q>q^{1/2}$. For every even $k$ choose
$e_k=2$. At every required degree $n=2r$ all counts vanish, so the
opposite-sign hypothesis holds with empty $B_k$ and zero error. The top
spectrum disappeared under powering. This refutes the earlier draft
extension without visibility, while leaving the positive-multiset
statement \fid{F17} intact. Geometrically the pointwise sequence itself
is the alternating trace of the constant line minus the unramified line
with Frobenius $-1$; its failure is not an exotic nonalgebraic sequence.

Second, the real-coefficient condition cannot simply be dropped for the
prescribed unbounded even set. Let $a=0$, $d=1$, and $t_n(x)=i$ for every
$x,n$. Then $M_n=iq^n$ has the root $q>q^{1/2}$ with coefficient $i$.
For $k=2m\ge6$ with $k\equiv2\pmod4$, put $a_j=x_{2j-1}$ and $b_j=x_{2j}$. The path
polynomial is
\[
 P_k=\sum_{j=1}^{m}a_jb_j+
       \sum_{j=1}^{m-1}a_{j+1}b_j.
\]
Its coefficient matrix is triangular with diagonal entries one, so it is
equivalent to a nondegenerate bilinear pairing. Its zero count over
$\F_Q$ is $Q^m+(Q^m-1)Q^{m-1}$, including characteristic two.
Since $i^k=-1$,
\[
 QN_{k,n}=-Q^k-Q^{k/2+1}+Q^{k/2}.
\]
Take $B_k=\{q^k,q^{k/2+1}\}$ and $E_{k,n}=Q^{k/2}$; the error has
$\eta_k=0$. Every other stated hypothesis holds for all $n$. This is a
counterexample to that deleted hypothesis, not to the integer-Euler
version or to possible other criteria with phase control.

\paragraph{What the stronger theorem removes.}
For a bounded constructible Weil complex, alternating cohomological
multiplicities are integers. The input need not be supported in one parity
if the even-path output has the prescribed negative visible spectrum and
the pointwise bound holds. Identical roots that already cancel between
degrees remain invisible: for example
$\Q_\ell(-1)\oplus\Q_\ell(-1)[-1]$ has zero trace in every extension.
Consequently \fid{F24} strictly weakens the input parity gate; it does not
infer bounds for erased cohomology or provide the path-output sign.

